%presentation
\documentclass{beamer}

%impressions
%\documentclass[handout]{beamer}
%\usepackage{pgfpages}
%\pgfpagesuselayout{2 on 1}[a4paper,border shrink=5mm]
%\setbeameroption{notes on second screen}
%\pgfpagelayout{2 on 1}[a4paper, border, shrink=5mm]
% vue sur http://wwwtaketorg/spip/articlephp3?id_article=30...
\usepackage[T1]{fontenc}
\usepackage[frenchb]{babel}
\usepackage[utf8x]{inputenc} % Pour pouvoir taper les accents sans faire de combinaison
%\usepackage{arev}
%\usepackage{aeguill}
%mode code
\usepackage{listings}

%mode verbatim
\usepackage{moreverb}

%\usepackage[darktab]{beamerthemesidebar}
%\leftsidebar
%\usetheme{Hannover}
%\usetheme{Warsaw}
%\usetheme{PaloAlto}
\usetheme{JuanLesPins}
%\usetheme{Antibes}
%\usetheme{Shingara}
%\usetheme{Berlin}
%\usetheme{Oxygen}
\usepackage{thumbpdf}
\usepackage{wasysym}
\usepackage{ucs}
\usepackage{pgfarrows,pgfnodes,pgfautomata,pgfheaps,pgfshade}
\usepackage{verbatim}
\usepackage{color}

\AtBeginSection[]
{
  \begin{frame}<beamer>
    \tableofcontents[currentsection,currentsubsection]
  \end{frame}
}


\title{Ruby On Rails par Rspec}
\author{Cyril Mougel}

\logo{\includegraphics[width=10mm]{rails.png}}

\begin{document}

\begin{frame}
    \titlepage
\end{frame}

\section{Présentation}

\begin{frame}
	\frametitle{Typo 5.1.3}
	\begin{itemize}
		\item supporte uniquement Rails 2.0.2
		\item Pas d'évolution vers Rails 2.1
		\item Couverture de code
	\end{itemize}
\end{frame}

\section{Migration lié à Rails 2.1}

\begin{frame}
    \frametitle{Mise à jour du projet}
    \begin{itemize}
        \item rake rails:update
        \item Mise à jour des fichiers de boot
        \item Mise à jour des fichiers de JS
    \end{itemize}
\end{frame}

\begin{frame}
    \frametitle{Plus de follow\_redirect dans les test de controlleur}
    \begin{itemize}
        \item possibilité d'utiliser follow\_redirect dans un Test::ActionController
        \item utilisation uniquement dans les test d'intégration
        \item inutile dans les test d'action
    \end{itemize}
\end{frame}

\begin{frame}
    \frametitle{Ce qui est déprécié ou supprimé}
    \begin{itemize}
        \item render\_partial est supprimé suite à sa déprécition sur Rails 2.0
    \end{itemize}
\end{frame}

\section{Migration lié à Rails 2.2}

\begin{frame}
  \frametitle{relative\_url\_root en configuration de ActionController}
  \begin{itemize}
    \item Le helper link\_to utilise relative\_url\_root
    \item Avant : utilisait @request.relative\_url\_root
    \item Maintenant : On défini directement dans la configuration :
    ActionController::Base.relative\_url\_root
  \end{itemize}
\end{frame}

\begin{frame}
    \frametitle{Migration dans des transactions}
    \begin{itemize}
        \item Si la migration échoue, pas de mise à jour de la BDD
        \item raise Exception == Migration Failed, même avec rescue
        \item Eviter tous les cas de Raise.
    \end{itemize}
\end{frame}

\begin{frame}
    \frametitle{Chargement des classes de controller avec cache\_classe = true}
    \begin{itemize}
        \item en production cache\_classe = true
        \item Chargement des classes en mémoire
        \item rake db:migrate RAILS\_ENV='production' charge classe 
        \item si appel ActiveRecord en cache du controller == FAILED
        \item rescue it
    \end{itemize}
\end{frame}

\begin{frame}
  \frametitle{et des trivialités}
  \begin{itemize}
    \item fragment\_cache\_store remplacer par cache\_store
    \item TextHelper::truncate prend un seul argument avec :length en option.
    La méthode avec plusieurs argument DEPRECATED
  \end{itemize}
\end{frame}


\end{document}
